\writeSectionFrame{\presentationTitle}{Beobachter}
\begin{frame}{Steckbrief}
	\begin{itemize}
 		\item Deutscher Name: Beobachter
 		\item Englischer Name: Observer
 		\item Auch bekannt als: Publisher-Subscriber
 		\item Gruppe: Verhaltensmuster
 	\end{itemize}
\end{frame}

\begin{frame}{Beobachter}
	\begin{block}{Zweck}
 		Definiere eine 1-zu-n-Abhängigkeit zwischen Objekten, so dass eine Änderung des Zustandes eines Objekts dazu führt, dass alle abhängigen Objekte benachrichtigt und automatisch aktualisiert werden.
 	\end{block}
\end{frame}

\frameWithImageOnly{\BILDERPFAD/observer1.jpg}

\begin{frame}{Allgemeines Klassendiagramm}
    \begin{tikzpicture}[scale=0.7, transform shape]
        \umlclass[x=0,y=0]{Subject}{%
            +addObserver(observer: Observer): void\\
            +removeObserver(observer: Observer): void\\
            +notifyObservers(): void
        }{}
    
        \umlclass[x=7,y=0]{ConcreteSubject}{%
            -observers: List<Observer>\\
            -state: String\\
            +setState(state: String): void\\
            +getState(): String
        }{}
    
        \umlclass[x=0,y=-5]{Observer}{%
            +update(): void
        }{}
    
        \umlclass[x=7,y=-5]{ConcreteObserver}{%
            -subject: ConcreteSubject\\
            +update(): void
        }{}
    
        % Define the relationships
        \umlinherit{ConcreteSubject}{Subject}
        \umlinherit{ConcreteObserver}{Observer}
    
        \umluniassoc{Subject}{Observer}
        \umluniassoc{ConcreteObserver}{ConcreteSubject}
    
    \end{tikzpicture}
\end{frame}

\begin{frame}{Beispiel News Channel}
    \begin{tikzpicture}[scale=0.7, transform shape]
    \umlclass[x=0,y=0]{NewsAgency}{%
        +addChannel(newsChannel: NewsChannel): void\\
        +removeChannel(newsChannel: NewsChannel): void\\
        +notifyChannels(): void
    }{}

    \umlclass[x=7,y=0]{MyNewsAgency}{%
    }{}

    \umlclass[x=0,y=-5]{NewsChannel}{%
        +update(): void
    }{}

    \umlclass[x=7,y=-3]{RadioNewsChannel}{%
    }{}

    \umlclass[x=7,y=-5]{TvNewsChannel}{%
    }{}

    \umlclass[x=7,y=-7]{OnlineNewsChannel}{%
    }{}

    % Define the relationships
    \umlinherit{MyNewsAgency}{NewsAgency}
    \umlinherit{RadioNewsChannel}{NewsChannel}
    \umlinherit{TvNewsChannel}{NewsChannel}
    \umlinherit{OnlineNewsChannel}{NewsChannel}

    \umluniassoc{NewsChannel}{NewsAgency}
    \umluniassoc{MyNewsAgency}{NewsChannel}

\end{tikzpicture}
\end{frame}

\begin{frame}[fragile]{Subjekt}
	\begin{exampleblock}{NewsAgency}
 		\begin{lstlisting}
public interface NewsAgency {

	void addChannel(NewsChannel channel);

	void removeChannel(NewsChannel channel);

	void setNews(String news);

}
 		\end{lstlisting}
 	\end{exampleblock}
\end{frame}

\begin{frame}[fragile]{Konkretes Subjekt}
	\begin{exampleblock}{NewsAgencyImpl}
 		\begin{lstlisting}
public class NewsAgencyImpl implements NewsAgency {
	private List<NewsChannel> channels;
	
	public NewsAgencyImpl() {
		channels = new ArrayList<>();
	}
	
	@Override
	public void addChannel(NewsChannel channel) {
		channels.add(channel);
	}
} 		    
 		\end{lstlisting}
 	\end{exampleblock}
\end{frame}

\begin{frame}[fragile]{Konkretes Subjekt}
	\begin{exampleblock}{NewsAgencyImpl}
 		\begin{lstlisting}
@Override
public void removeChannel(NewsChannel channel) {
    channels.remove(channel);
}

@Override
public void setNews(String news) {
    for (NewsChannel channel: channels) {
        channel.update(news);
    }
}
 		\end{lstlisting}
 	\end{exampleblock}
\end{frame}

\begin{frame}[fragile]{Beobachter}
	\begin{exampleblock}{NewsChannel}
 		\begin{lstlisting}
public interface NewsChannel {
	void update(String object);
}
 		\end{lstlisting}
 	\end{exampleblock}
\end{frame}

\begin{frame}[fragile]{Konkreter Beobachter}
	\begin{exampleblock}{TvNewsChannel}
 		\begin{lstlisting}
public class TvNewsChannel implements NewsChannel {

	@Override
	public void update(String news) {
		System.out.println("Im Fernsehen laufen die "
				+ "neuesten Nachrichten");
		System.out.println(news);
	}
} 		    
 		\end{lstlisting}
 	\end{exampleblock}
\end{frame}

\frameWithImageOnly{\BILDERPFAD/observer3.jpg}
\note{
    Im folgenden Beispiel wird eine kleine Kalender-App simuliert. Dabei können sich Beobachter für bestimmte Operationen anmelden, d.h. ein Beobachter möchte nicht immer informiert werden, sondern nur wenn eine bestimmte Operation auf dem Kalender ausgeführt wird, also z.B. neuen Termin erstellen, Termin löschen usw. Außerdem soll das Subjekt den Beobachtern weitere Informationen beim Aktualisieren mitgeben.
}

\begin{frame}{Beispiel Kalender}
    \begin{tikzpicture}[scale=0.7, transform shape]
    \umlclass[x=-4,y=0]{Subject}{%
    }{
        void addObserver(observer: Observer , dateOperation: DateOperation)\\
        void removeObserver(observer: Observer , dateOperation: DateOperation)\\
        void updateObservers(dateOperation: DateOperation, furtherInformation: Object)\\
    }

    \umlclass[x=1,y=-3]{DateChooser}{%
    }{
        void chooseDate(date: Date, dateOperation: DateOperation)
    }

    \umlclass[x=-5,y=-6]{Observer}{%
    }{
        update(dateOperation: DateOperation, furtherInformation: Object);
    }

    \umlclass[x=5,y=-6]{Calendar}{%
    }{}

    % Define the relationships
    \umlinherit{DateChooser}{Subject}
    \umlinherit{Calendar}{Observer}

    \umluniassoc{Observer}{Subject}
    \umluniassoc{DateChooser}{Calendar}

\end{tikzpicture}
\end{frame}

\begin{frame}[fragile]{Subjekt}
	\begin{exampleblock}{Subject}
 		\begin{lstlisting}
public abstract class Subject {
	private Map<DateOperation, List<Observer>> observers;
	
	public Subject() {
		observers = new HashMap<>();
	}
	
	public void addObserver(Observer observer, 
			DateOperation dateOperation) {
		List<Observer> observerList = observers.get(dateOperation);
		
		if (observerList == null) {
			observerList = new ArrayList<>();
		}
		
		observerList.add(observer);
		observers.put(dateOperation, observerList);
	}
 		\end{lstlisting}
 	\end{exampleblock}
\end{frame}

\begin{frame}[fragile]{Subjekt}
	\begin{exampleblock}{Subject}
 		\begin{lstlisting}
public void removeObserver(Observer observer, 
        DateOperation dateOperation) {
    List<Observer> observerList = observers.get(dateOperation);
    observerList.remove(observer);
}

public void updateObservers(DateOperation dateOperation,
        Object object) {
    List<Observer> observerList = observers.get(dateOperation);
    for (Observer observer: observerList) {
        observer.update(dateOperation, object);
    }
}
 		\end{lstlisting}
 	\end{exampleblock}
\end{frame}

\begin{frame}[fragile]{Konkretes Subjekt}
	\begin{exampleblock}{DateChooser}
 		\begin{lstlisting}
public class DateChooser extends Subject {
	private Date date;
	
	public void chooseDate(Date date, 
            DateOperation dateOperation) {
		this.date = date;
		updateObservers(dateOperation, date);
	}
	
	public Date getDate() {
		return date;
	}
 }
 		\end{lstlisting}
 	\end{exampleblock}
\end{frame}

\begin{frame}[fragile]{Beobachter}
	\begin{exampleblock}{Observer}
 		\begin{lstlisting}
public interface Observer {
	void update(DateOperation dateOperation, 
			Object furtherInformation);
}

 		\end{lstlisting}
 	\end{exampleblock}
\end{frame}

\begin{frame}[fragile]{Konkreter Beobachter}
	\begin{exampleblock}{Calendar}
 		\begin{lstlisting}
public class Calendar implements Observer {
	private final Map<Date, String> appointments;
	private final DateChooser dateChooser;

	@Override
	public void update(DateOperation dateOperation, 
			Object furtherInformation) {
		switch(dateOperation) {
		case ADD_DATE: // ...
		break;
		case REMOVE_DATE: // ...
		}
	}
} 		    
 		\end{lstlisting}
 	\end{exampleblock}
\end{frame}

\frameWithImageOnly{\BILDERPFAD/observer4.jpg}
\note{
    Das nächste Beispiel simuliert eine Wetterstation. Dabei sollen die Wetterdaten auf verschiedene Arten angezeigt werden. Einmal haben wir eine Darstellung der aktuellen Wetterdaten und dann gibt es noch eine statistische Darstellung. Dieses Beispiel wäre auch wunderbar als verteiltes System denkbar.
}

\begin{frame}{Beispiel Wetterstation}
    \begin{tikzpicture}[scale=0.7, transform shape]
        \umlclass[x=0,y=0]{WeatherData}{%
        }{}
    
        \umlclass[x=7,y=0]{WeatherStation}{%
        }{}
    
        \umlclass[x=0,y=-5]{WeatherObserver}{%
        }{}
    
        \umlclass[x=6,y=-3]{CurrentConditionsDisplay}{
        }{}

        \umlclass[x=11,y=-5]{StatisticsDisplay}{%
        }{}
    
        % Define the relationships
        \umlinherit{WeatherStation}{WeatherData}
        \umlinherit{CurrentConditionsDisplay}{WeatherObserver}
        \umlinherit{StatisticsDisplay}{WeatherObserver}
    
        \umluniassoc{WeatherData}{WeatherObserver}
        \umluniassoc{CurrentConditionsDisplay}{WeatherStation}
        \umluniassoc{StatisticsDisplay}{WeatherStation}
    
    \end{tikzpicture}
\end{frame}

\begin{frame}[fragile]{Subjekt}
	\begin{exampleblock}{Subject}
 		\begin{lstlisting}
public interface WeatherData {
    void addObserver(WeatherObserver observer);
    void removeObserver(WeatherObserver observer);
    void notifyObservers();
}
 		\end{lstlisting}
 	\end{exampleblock}
\end{frame}

\begin{frame}[fragile]{Konkretes Subjekt}
	\begin{exampleblock}{Subject}
 		\begin{lstlisting}
public class WeatherStation implements WeatherData {
    private List<WeatherObserver> observers;

    @Override
    public void addObserver(WeatherObserver observer) {
        observers.add(observer);
    }

    @Override
    public void notifyObservers() {
        for (WeatherObserver observer : observers) {
            observer.update(temperature, humidity, pressure);
        }
    }
}
 		\end{lstlisting}
 	\end{exampleblock}
\end{frame}

\begin{frame}[fragile]{Beobachter}
	\begin{exampleblock}{Observer}
 		\begin{lstlisting}
public interface WeatherObserver {
    void update(float temperature, 
                float humidity, 
                float pressure);
}
 		\end{lstlisting}
 	\end{exampleblock}
\end{frame}

\begin{frame}[fragile]{Konkreter Beobachter}
	\begin{exampleblock}{Calendar}
 		\begin{lstlisting}
public class CurrentConditionsDisplay 
        implements WeatherObserver {
    private float temperature;
    private float humidity;

    @Override
    public void update(float temperature, 
            float humidity, float pressure) {
        this.temperature = temperature;
        this.humidity = humidity;
        display();
    }
}	    
 		\end{lstlisting}
 	\end{exampleblock}
\end{frame}